\documentclass[11pt]{article}
\usepackage[a4paper,margin=25mm]{geometry}
\usepackage{amsmath,amssymb,amsthm,hyperref}
\usepackage{microtype}\emergencystretch=2em
\newtheorem{theorem}{Theorem}\newtheorem{lemma}[theorem]{Lemma}\newtheorem{corollary}[theorem]{Corollary}
\newcommand{\GL}{\operatorname{GL}}\newcommand{\SL}{\operatorname{SL}}
\newcommand{\Sym}{\operatorname{Sym}}\newcommand{\im}{\operatorname{im}}
\title{Root parameters and fixed-space obstructions}
\author{A source-attributed recovery for the Markovskaya--Nuzhin classifications}
\date{}
\begin{document}\maketitle
This common paper supplies elementary proofs of the root-recognition and
fixed-space tools used in the accompanying linear-group application. The results
are classical; no priority claim is made. In particular, the root recognition
specializes Levchuk's theorem, and the fixed-space inequality is Scott's.
All fields in the root-recognition statements are finite. Write
$t_{ij}(a)=I+aE_{ij}$, and use $x^y=yxy^{-1}$ and
$[x,y]=xyx^{-1}y^{-1}$.

\begin{lemma}[Root propagation]\label{propagate}
Let $M\leq\GL_n(F)$ and let a directed graph on $\{1,\ldots,n\}$ have an
arc $i\to j$ when $M$ contains some $t_{ij}(a)$ with $a\ne0$.
If this graph is strongly connected, $M$ meets every root subgroup nontrivially.
In particular a monomial matrix with permutation part an $n$-cycle, together
with a nonidentity root element between cyclically consecutive positions,
has this property, for either orientation.
\end{lemma}
\begin{proof}
For distinct $i,j,k$, elementary matrix multiplication gives
$[t_{ij}(a),t_{jk}(b)]=t_{ik}(ab)$. A simple directed path therefore supplies
its endpoint root by induction; its vertices are distinct, so every use of
this formula is legitimate. Conjugation by a monomial matrix transports a
root to the root on the permuted ordered pair and multiplies its parameter
by a nonzero scalar. The conjugates of the indicated adjacent root give a
directed Hamiltonian cycle. The assertion for $n=2$ needs only these two
conjugates, not a three-index commutator.
\end{proof}

\begin{theorem}[Finite root-parameter recognition]\label{recognition}
Let $n\geq3$ and $M\leq\GL_n(F_q)$ meet every root subgroup nontrivially.
Put $A_{ij}=\{a:t_{ij}(a)\in M\}$ and
$R=\langle t_{ij}(A_{ij}):i\ne j\rangle$.
There are a subfield $K\leq F_q$ and an invertible diagonal matrix $D$ such
that $DRD^{-1}=\SL_n(K)$. This does not assert that $DMD^{-1}$ equals
$\SL_n(K)$.
Consequently every determinant-one diagonal sign matrix belongs to $R$,
and $R$ contains a monomial lift of every coordinate permutation. If
$a,b\in A_{ij}$ with $b\ne0$, then $a/b\in K$; if
$a\in A_{ij}$ and $b\in A_{ji}$, then $ab\in K$.
\end{theorem}
\begin{proof}
Each $A_{ij}$ is a nonzero finite additive group and
$A_{ij}A_{jk}\subseteq A_{ik}$ for distinct indices. Multiplication by a
fixed nonzero element gives $|A_{ij}|\leq|A_{ik}|$; interchanging $j,k$
gives equality. Thus the size is constant in each row. Similarly left
multiplication shows it is constant in each column. For $n\geq3$ the
off-diagonal positions are connected by row and column steps, so all sizes
are the same number $m$. Every preceding product with one fixed nonzero
factor is therefore an equality.

Choose $a_i\in A_{i1}\setminus\{0\}$ for $i>1$ and conjugate by
$D=\operatorname{diag}(1,a_2^{-1},\ldots,a_n^{-1})$. Write $B_{ij}$ for
the new parameter groups. Then $1\in B_{i1}$. For distinct $i,j>1$,
$B_{ij}\subseteq B_{i1}$, by multiplying by $1\in B_{j1}$, and equality
follows from the sizes. Also $B_{1j}\subseteq B_{ij}$ by multiplying by
$1\in B_{i1}$, so
\[
 B_{ij}=B_{i1}=B_{1j}=:K_i\qquad(i,j>1,\ i\ne j).
\]
The product $B_{i1}B_{1j}\subseteq B_{ij}$ says $K_iK_i\subseteq K_i$.
Thus $K_i$ is a finite subring of the field containing $1$, hence a field:
multiplication by a nonzero element is injective and surjective on $K_i$.
Finally $B_{ij}B_{j1}\subseteq B_{i1}$ gives $K_iK_j\subseteq K_i$, hence
$K_j\subseteq K_i$ and equality by sizes. There is at least one such $j$
for each $i$ because $n\geq3$. All the $B_{ij}$ are consequently the same
field $K$. Elementary row elimination proves that their transvections
generate $\SL_n(K)$.

Diagonal conjugation preserves diagonal matrices and monomial permutation
parts. Determinant-one sign matrices lie in $\SL_n(K)$, as do monomial
lifts of transpositions (use the $2\times2$ block
$\left(\begin{smallmatrix}0&-1\\1&0\end{smallmatrix}\right)$).
The ratio assertion and the opposite-root product assertion follow because
the diagonal conjugating factors cancel.
\end{proof}

\begin{corollary}[Primitive-square test and isolation]\label{square}
Suppose $q$ is odd and $u$ generates $F_q^*$. Then $F_p(u^2)=F_q$.
Accordingly, in Theorem~\ref{recognition}, either two same-root parameters
with ratio $u^2$, or opposite-root parameters with product a nonzero
prime-field multiple of $u^2$, force $R=\SL_n(F_q)$.
For $n\geq4$, if $i,j,k$ are distinct and
$t_{ij}(a)t_{kj}(b)\in M$, each of these two root factors can be isolated
in $M$ when $q$ is odd.
\end{corollary}
\begin{proof}
If $u^2$ lay in a proper subfield of size $p^d$, then $p^d\leq\sqrt q$
and $(q-1)/2\leq p^d-1$, impossible since $q+1>2\sqrt q$. A prime field
has no proper subfield, covering that case as well.
For isolation choose an index $\ell$ different from $i,j,k$. The diagonal
matrix with minus signs at $k,\ell$ belongs to $R$. Its commutator with
the displayed product is $t_{kj}(-2b)$. An integer power multiplies this
coefficient by $(-2)^{-1}\in F_p$, yielding $t_{kj}(b)$. Divide it out to
obtain the other factor. This uses the previously established $R$; it is
not a substitute for proving that all roots were reached.
\end{proof}

\begin{lemma}[Symmetric square]\label{symmetric}
For a field $F$ of characteristic different from $2$ and $n\geq2$, the
natural action of $\GL_n(F)$ on $\Sym^2(F^n)$ is irreducible and has
kernel $\{I,-I\}$.
\end{lemma}
\begin{proof}
Use the basis $e_i^2,e_ie_j$ ($i<j$). The diagonal sign group has distinct
characters on the off-diagonal basis vectors, and the trivial character
on the span of the squares. Averaging against its characters is permitted,
since its order is a power of $2$. Hence an invariant subspace with any
nonzero off-diagonal coefficient contains an individual $e_ie_j$.
Permutation matrices give all such vectors. Replacing $e_i$ by $e_i+e_j$
then yields $e_j^2$ by subtraction, and all squares follow. If instead a
nonzero vector has only square coefficients, choose $a_i\ne0$ in
$\sum a_ie_i^2$ and replace $e_i$ by $e_i+e_j$; the new off-diagonal
coefficient is $2a_i\ne0$. Thus every nonzero invariant subspace is full.
If $g$ acts trivially, $(ge_i)^2=e_i^2$ forces $ge_i=\epsilon_i e_i$,
$\epsilon_i\in\{1,-1\}$. The mixed products force all $\epsilon_i$ equal.
Both scalar signs do act trivially.
\end{proof}

\begin{theorem}[Scott's inequality, with proof]\label{scott}
Suppose $g_1\cdots g_m=1$ generate an irreducible subgroup of $\GL(V)$,
where $\dim V\geq2$. Then
\[\sum_i\operatorname{codim}\operatorname{Fix}(g_i)\geq2\dim V.\]
\end{theorem}
\begin{proof}
Put $p_i=g_1\cdots g_{i-1}$ and
$W=\bigoplus_i\im(g_i-I)$. The map
$\Phi:W\to V$, $(w_i)\mapsto\sum_i p_iw_i$, is onto.
Indeed a covector annihilating its image annihilates $g_1-I$ and hence is
fixed by $g_1$; inductively it is fixed by every $g_i$. The dual has no
nonzero common fixed vector, since its annihilator would be a proper
invariant hyperplane in $V$.
The injective map $v\mapsto((g_i-I)v)_i$ lands in $\ker\Phi$, by the
telescoping identity $\sum_i p_i(g_i-I)=g_1\cdots g_m-I=0$.
Its kernel is the common fixed space, zero by irreducibility and
$\dim V>1$. Therefore $\dim W\geq\dim V+\dim V$.
\end{proof}

\begin{corollary}\label{gl43}
$\GL_4(3)$ cannot be generated by three involutions two of which commute.
\end{corollary}
\begin{proof}
If $a,b,c$ were such generators with $ab=ba$, apply the preceding results
to the ten-dimensional symmetric square and the five factors
$a,b,ab,c,c$, whose product is $1$. An involution with an $r$-dimensional
negative eigenspace has symmetric-square fixed-space codimension
$r(4-r)\leq4$. Scott's lower bound is $20$, so all five codimensions must
be $4$, in particular those of $a,b,c$. They all have $r=2$, determinant
$1$, and cannot generate the full general linear group.
\end{proof}

\paragraph{Attribution and boundary.}
The companion source is I.~A.~Markovskaya and Ya.~N.~Nuzhin,
\emph{On generation of the groups $\GL_n(q)$ and $\operatorname{PGL}_n(q)$
by three involutions, two of which commute}, Trudy Inst. Mat. Mekh. UrO RAN
31(4) (2025), 247--259, DOI 10.21538/0134-4889-2025-31-4-fon-03.
Its Lemmas 6--9 invoke Nuzhin's symmetric-square theorem, Scott's
\emph{Matrices and cohomology} (1977), and Levchuk's root-generation
theorem (1983). The elementary proofs here replace those particular
invocations, not their general published theorems. No finite search,
symbolic software, or formal certificate is a premise of this paper.
\end{document}
